\section{Introduction}\label{sec:intro}

For a prime-order subgroup $\langle G\rangle$ of size $r$ on an elliptic
curve $E/\FF_q$, the best generic algorithm for the discrete-logarithm
problem is Pollard rho~\cite{pollard1978,vanoorschot1999}, with expected
cost $\sqrt{\pi r/(2A)}$ group operations when $A$ automorphisms (negation,
and on a Koblitz curve the Frobenius) can be folded into the
walk~\cite{wiener1998,gallant2000}. Index calculus, the method that breaks
finite-field discrete logarithms in subexponential time, was brought to
elliptic curves by Semaev's summation polynomials~\cite{semaev2004} and by
Gaudry's and Diem's analyses over extension
fields~\cite{gaudry2009,diem2011}. Its open question has always been
\emph{where the crossover is}: on which curves, at which sizes and under
which accounting does it cost less than rho?

This paper documents fourteen weeks (1 July--14 October 2026) of work on
that question in one repository. The index-calculus programme proper ran
from 6 September to 14 October; July and August hold the preceding
prime-field structure work that motivated it
(\S\ref{sec:chronology}). The output is large: about 6{,}900 commits,
about 450 dated study directories under \art{research/}, 27 rows of the
performance index \art{benchmarks/README.md}, a boundary ledger
(\art{docs/ic/boundary_targets.json}), a native measurement harness
(\art{src/cryptanalysis/ecbench/}), and a companion manuscript on
root-finding decomposition (\art{papers/rootfinding-decomposition/}). It
is easy, in a body of work this size, to improve a headline number by two
orders of magnitude and establish nothing; this paper is organised to
prevent that. Every improvement is stated with the boundary on which it
was measured (oracle stage, complete solve, cold one-target cost, online
cost after retained set-up), the unit, the sample, the interval, and the
committed artefact that holds the records.

\paragraph{The routes.} We pursued six lines, each a section below:
\begin{enumerate}
\item \emph{Root-finding decomposition} (\S\ref{sec:rootfinding}):
  replacing the Gr\"obner solve of the decomposition step by closed-form
  roots of $\Sem3$, then by meet-in-the-middle tables for more summands,
  then by GPU, tower and dense-index variants.
\item \emph{Gr\"obner-basis engines} (\S\ref{sec:groebner}): matrix F4,
  F5, an inherited F4 that re-uses its root echelon form, an
  index-calculus-specific F6-IC, an $\FF_p$ F4 brought to msolve parity,
  and an audit that found and fixed one soundness bug.
\item \emph{Koblitz and binary curves} (\S\ref{sec:koblitz}):
  Frobenius-invariant factor bases~\cite{ggmp2020}, their compact
  signed-orbit form, SAT oracles~\cite{trimoska2020}, and a ladder of
  one-target online comparisons with rho from $n=53$ to $n=83$.
\item \emph{Prime and extension fields} (\S\ref{sec:prime}): $j=0$ and
  $a=-3$ known-answer ladders, the Kummer-tower factor base of
  Petit--Kosters--Messeng~\cite{pkm2016}, and Gaudry's $k=4$ index
  calculus over $\FF_{p^4}$ with the $D_4$ symmetry of
  Faug\`ere--Gaudry--Huot--Renault~\cite{fghr2014}.
\item \emph{Isogeny, transfer and structure routes}
  (\S\ref{sec:transfer}): conductor gaps in the ECC2K-130 isogeny class,
  isogeny fibres, transfer-aware index calculus, GHS and Weil descent;
  mostly negative, with one lemma that explains why.
\item \emph{Literature replication} (\S\ref{sec:prime},
  \S\ref{sec:transfer}): the \emph{lit2026} programme, which turned each
  finding of a 2025--26 literature search into a measured change, a
  diagnostic, or a recorded negative.
\end{enumerate}
Underneath all six sits the measurement work of \S\ref{sec:method}: a
single harness (\emph{ecbench}) that charges rho, BSGS, the kangaroo and
index calculus in one unit on one-target workloads, and the accounting
corrections that the harness forced on earlier claims.

\paragraph{Contributions.}
\begin{itemize}
\item A unit and a boundary: $S=\text{group-addition
  equivalents}/\sqrt{r}$, cold one-target with every phase charged, beside
  a separately reported online window; sealed, replayable records; and a
  frontier definition under which two numbers are comparable only when
  every component of the measurement domain matches
  (\S\ref{sec:method}).
\item Root-finding decomposition taken from an observation to an
  engineered, verified, harness-integrated oracle family, with a
  machine-checked proof of the identity it rests on
  (\S\ref{sec:rootfinding}).
\item A documented map of which levers move a constant, which move nothing,
  and which make things worse, across about 40 measured variants of the
  decomposition, factor-base and linear-algebra stages
  (Table~\ref{tab:improvements}).
\item An honest frontier (\S\ref{sec:frontier}): the cold one-target
  frontier is still rho; the retained-set-up online frontier has an
  admitted index-calculus point; the stage frontier moved substantially.
\end{itemize}

\paragraph{What this paper does not claim.} No asymptotic sub-rho
algorithm, no change to the Gaudry--Diem picture for prime fields, no
statement about ECC2K-130, P-256 or any deployed curve beyond labelled
arithmetic on public parameters, no key recovery. Every instance is a
public synthetic or standard-parameter known-answer instance. Every change
is classed, in the repository's vocabulary, as an \cls{advance} (an
exponent or a crossover moves under complete accounting),
\cls{engineering} (a constant moves), \cls{relabelling}, or
\cls{accounting}; with one modelled exception (\S\ref{sec:gaudry}) every
improvement in this paper is \cls{engineering} or \cls{accounting}.

\paragraph{How to read the numbers.} ``$X\times$'' always means
baseline cost divided by candidate cost on the stated boundary. Stage
ratios (one decomposition, one Gr\"obner call) are never presented as
end-to-end speed-ups; wall-time ratios from hosts that did not earn an
isolation level are labelled \emph{descriptive}; a cost with unpriced work
is labelled a \emph{lower bound}. Appendix~\ref{app:artefacts} lists the
artefact behind each table.
